\IfFileExists{stacks-project.cls}{%
\documentclass{stacks-project}
}{%
\documentclass{amsart}
}

% For dealing with references we use the comment environment
\usepackage{verbatim}
\newenvironment{reference}{\comment}{\endcomment}
%\newenvironment{reference}{}{}
\newenvironment{slogan}{\comment}{\endcomment}
\newenvironment{history}{\comment}{\endcomment}

% For commutative diagrams we use Xy-pic
\usepackage[all]{xy}

% We use 2cell for 2-commutative diagrams.
\xyoption{2cell}
\UseAllTwocells

% We use multicol for the list of chapters between chapters
\usepackage{multicol}

% This is generally recommended for better output
\usepackage{lmodern}
\usepackage[T1]{fontenc}

% For cross-file-references
\usepackage{xr-hyper}

% Package for hypertext links:
\usepackage{hyperref}

% For any local file, say "hello.tex" you want to link to please
% use \externaldocument[hello-]{hello}
\externaldocument[introduction-]{introduction}
\externaldocument[conventions-]{conventions}
\externaldocument[sets-]{sets}
\externaldocument[categories-]{categories}
\externaldocument[topology-]{topology}
\externaldocument[sheaves-]{sheaves}
\externaldocument[sites-]{sites}
\externaldocument[stacks-]{stacks}
\externaldocument[fields-]{fields}
\externaldocument[algebra-]{algebra}
\externaldocument[brauer-]{brauer}
\externaldocument[homology-]{homology}
\externaldocument[derived-]{derived}
\externaldocument[simplicial-]{simplicial}
\externaldocument[more-algebra-]{more-algebra}
\externaldocument[smoothing-]{smoothing}
\externaldocument[modules-]{modules}
\externaldocument[sites-modules-]{sites-modules}
\externaldocument[injectives-]{injectives}
\externaldocument[cohomology-]{cohomology}
\externaldocument[sites-cohomology-]{sites-cohomology}
\externaldocument[dga-]{dga}
\externaldocument[dpa-]{dpa}
\externaldocument[sdga-]{sdga}
\externaldocument[hypercovering-]{hypercovering}
\externaldocument[schemes-]{schemes}
\externaldocument[constructions-]{constructions}
\externaldocument[properties-]{properties}
\externaldocument[morphisms-]{morphisms}
\externaldocument[coherent-]{coherent}
\externaldocument[divisors-]{divisors}
\externaldocument[limits-]{limits}
\externaldocument[varieties-]{varieties}
\externaldocument[topologies-]{topologies}
\externaldocument[descent-]{descent}
\externaldocument[perfect-]{perfect}
\externaldocument[more-morphisms-]{more-morphisms}
\externaldocument[flat-]{flat}
\externaldocument[groupoids-]{groupoids}
\externaldocument[more-groupoids-]{more-groupoids}
\externaldocument[etale-]{etale}
\externaldocument[chow-]{chow}
\externaldocument[intersection-]{intersection}
\externaldocument[pic-]{pic}
\externaldocument[weil-]{weil}
\externaldocument[adequate-]{adequate}
\externaldocument[dualizing-]{dualizing}
\externaldocument[duality-]{duality}
\externaldocument[discriminant-]{discriminant}
\externaldocument[derham-]{derham}
\externaldocument[local-cohomology-]{local-cohomology}
\externaldocument[algebraization-]{algebraization}
\externaldocument[curves-]{curves}
\externaldocument[resolve-]{resolve}
\externaldocument[models-]{models}
\externaldocument[functors-]{functors}
\externaldocument[equiv-]{equiv}
\externaldocument[pione-]{pione}
\externaldocument[etale-cohomology-]{etale-cohomology}
\externaldocument[proetale-]{proetale}
\externaldocument[relative-cycles-]{relative-cycles}
\externaldocument[more-etale-]{more-etale}
\externaldocument[trace-]{trace}
\externaldocument[crystalline-]{crystalline}
\externaldocument[spaces-]{spaces}
\externaldocument[spaces-properties-]{spaces-properties}
\externaldocument[spaces-morphisms-]{spaces-morphisms}
\externaldocument[decent-spaces-]{decent-spaces}
\externaldocument[spaces-cohomology-]{spaces-cohomology}
\externaldocument[spaces-limits-]{spaces-limits}
\externaldocument[spaces-divisors-]{spaces-divisors}
\externaldocument[spaces-over-fields-]{spaces-over-fields}
\externaldocument[spaces-topologies-]{spaces-topologies}
\externaldocument[spaces-descent-]{spaces-descent}
\externaldocument[spaces-perfect-]{spaces-perfect}
\externaldocument[spaces-more-morphisms-]{spaces-more-morphisms}
\externaldocument[spaces-flat-]{spaces-flat}
\externaldocument[spaces-groupoids-]{spaces-groupoids}
\externaldocument[spaces-more-groupoids-]{spaces-more-groupoids}
\externaldocument[bootstrap-]{bootstrap}
\externaldocument[spaces-pushouts-]{spaces-pushouts}
\externaldocument[spaces-chow-]{spaces-chow}
\externaldocument[groupoids-quotients-]{groupoids-quotients}
\externaldocument[spaces-more-cohomology-]{spaces-more-cohomology}
\externaldocument[spaces-simplicial-]{spaces-simplicial}
\externaldocument[spaces-duality-]{spaces-duality}
\externaldocument[formal-spaces-]{formal-spaces}
\externaldocument[restricted-]{restricted}
\externaldocument[spaces-resolve-]{spaces-resolve}
\externaldocument[formal-defos-]{formal-defos}
\externaldocument[defos-]{defos}
\externaldocument[cotangent-]{cotangent}
\externaldocument[examples-defos-]{examples-defos}
\externaldocument[algebraic-]{algebraic}
\externaldocument[examples-stacks-]{examples-stacks}
\externaldocument[stacks-sheaves-]{stacks-sheaves}
\externaldocument[criteria-]{criteria}
\externaldocument[artin-]{artin}
\externaldocument[quot-]{quot}
\externaldocument[stacks-properties-]{stacks-properties}
\externaldocument[stacks-morphisms-]{stacks-morphisms}
\externaldocument[stacks-limits-]{stacks-limits}
\externaldocument[stacks-cohomology-]{stacks-cohomology}
\externaldocument[stacks-perfect-]{stacks-perfect}
\externaldocument[stacks-introduction-]{stacks-introduction}
\externaldocument[stacks-more-morphisms-]{stacks-more-morphisms}
\externaldocument[stacks-geometry-]{stacks-geometry}
\externaldocument[moduli-]{moduli}
\externaldocument[moduli-curves-]{moduli-curves}
\externaldocument[examples-]{examples}
\externaldocument[exercises-]{exercises}
\externaldocument[guide-]{guide}
\externaldocument[desirables-]{desirables}
\externaldocument[coding-]{coding}
\externaldocument[obsolete-]{obsolete}
\externaldocument[fdl-]{fdl}
\externaldocument[index-]{index}

% Theorem environments.
%
\theoremstyle{plain}
\newtheorem{theorem}[subsection]{Theorem}
\newtheorem{proposition}[subsection]{Proposition}
\newtheorem{lemma}[subsection]{Lemma}

\theoremstyle{definition}
\newtheorem{definition}[subsection]{Definition}
\newtheorem{example}[subsection]{Example}
\newtheorem{exercise}[subsection]{Exercise}
\newtheorem{situation}[subsection]{Situation}

\theoremstyle{remark}
\newtheorem{remark}[subsection]{Remark}
\newtheorem{remarks}[subsection]{Remarks}

\numberwithin{equation}{subsection}

% Macros
%
\def\lim{\mathop{\mathrm{lim}}\nolimits}
\def\colim{\mathop{\mathrm{colim}}\nolimits}
\def\Spec{\mathop{\mathrm{Spec}}}
\def\Hom{\mathop{\mathrm{Hom}}\nolimits}
\def\Ext{\mathop{\mathrm{Ext}}\nolimits}
\def\SheafHom{\mathop{\mathcal{H}\!\mathit{om}}\nolimits}
\def\SheafExt{\mathop{\mathcal{E}\!\mathit{xt}}\nolimits}
\def\Sch{\mathit{Sch}}
\def\Mor{\mathop{\mathrm{Mor}}\nolimits}
\def\Ob{\mathop{\mathrm{Ob}}\nolimits}
\def\Sh{\mathop{\mathit{Sh}}\nolimits}
\def\NL{\mathop{N\!L}\nolimits}
\def\CH{\mathop{\mathrm{CH}}\nolimits}
\def\proetale{{pro\text{-}\acute{e}tale}}
\def\etale{{\acute{e}tale}}
\def\QCoh{\mathit{QCoh}}
\def\Ker{\mathop{\mathrm{Ker}}}
\def\Im{\mathop{\mathrm{Im}}}
\def\Coker{\mathop{\mathrm{Coker}}}
\def\Coim{\mathop{\mathrm{Coim}}}

% Boxtimes
%
\DeclareMathSymbol{\boxtimes}{\mathbin}{AMSa}{"02}

%
% Macros for moduli stacks/spaces
%
\def\QCohstack{\mathcal{QC}\!\mathit{oh}}
\def\Cohstack{\mathcal{C}\!\mathit{oh}}
\def\Spacesstack{\mathcal{S}\!\mathit{paces}}
\def\Quotfunctor{\mathrm{Quot}}
\def\Hilbfunctor{\mathrm{Hilb}}
\def\Curvesstack{\mathcal{C}\!\mathit{urves}}
\def\Polarizedstack{\mathcal{P}\!\mathit{olarized}}
\def\Complexesstack{\mathcal{C}\!\mathit{omplexes}}
% \Pic is the operator that assigns to X its picard group, usage \Pic(X)
% \Picardstack_{X/B} denotes the Picard stack of X over B
% \Picardfunctor_{X/B} denotes the Picard functor of X over B
\def\Pic{\mathop{\mathrm{Pic}}\nolimits}
\def\Picardstack{\mathcal{P}\!\mathit{ic}}
\def\Picardfunctor{\mathrm{Pic}}
\def\Deformationcategory{\mathcal{D}\!\mathit{ef}}
\usepackage{cleveref}
\crefname{originalreference}{original reference}{original references}
\Crefname{originalreference}{Original reference}{Original references}
\makeatletter
\newlabel{categories-lemma-nonempty-limit}{{086J}{source}{}{sourcepointer-1}{}}
\newlabel{categories-lemma-nonempty-limit@cref}{{[originalreference][0][]086J}{[1][0][]source}{}{}{}}
\makeatother
\crefname{originalreference}{original reference}{original references}
\Crefname{originalreference}{Original reference}{Original references}
\makeatletter
\newlabel{decent-spaces-definition-very-reasonable}{{03I8}{source}{}{sourcepointer-2}{}}
\newlabel{decent-spaces-definition-very-reasonable@cref}{{[originalreference][0][]03I8}{[1][0][]source}{}{}{}}
\makeatother
\crefname{originalreference}{original reference}{original references}
\Crefname{originalreference}{Original reference}{Original references}
\makeatletter
\newlabel{decent-spaces-lemma-locally-separated-decent}{{088J}{source}{}{sourcepointer-3}{}}
\newlabel{decent-spaces-lemma-locally-separated-decent@cref}{{[originalreference][0][]088J}{[1][0][]source}{}{}{}}
\makeatother
\crefname{originalreference}{original reference}{original references}
\Crefname{originalreference}{Original reference}{Original references}
\makeatletter
\newlabel{decent-spaces-section-reasonable-decent}{{03I7}{source}{}{sourcepointer-4}{}}
\newlabel{decent-spaces-section-reasonable-decent@cref}{{[originalreference][0][]03I7}{[1][0][]source}{}{}{}}
\makeatother
\crefname{originallemma}{original lemma}{original lemmas}
\Crefname{originallemma}{Original lemma}{Original lemmas}
\makeatletter
\newlabel{lemma-directed-inverse-system-has-limit}{{07SF}{source}{}{sourcepointer-5}{}}
\newlabel{lemma-directed-inverse-system-has-limit@cref}{{[originallemma][0][]07SF}{[1][0][]source}{}{}{}}
\makeatother
\crefname{originalreference}{original reference}{original references}
\Crefname{originalreference}{Original reference}{Original references}
\makeatletter
\newlabel{limits-lemma-topology-limit}{{086V}{source}{}{sourcepointer-6}{}}
\newlabel{limits-lemma-topology-limit@cref}{{[originalreference][0][]086V}{[1][0][]source}{}{}{}}
\makeatother
\crefname{originalreference}{original reference}{original references}
\Crefname{originalreference}{Original reference}{Original references}
\makeatletter
\newlabel{spaces-properties-lemma-finite-fibres-presentation}{{03IJ}{source}{}{sourcepointer-7}{}}
\newlabel{spaces-properties-lemma-finite-fibres-presentation@cref}{{[originalreference][0][]03IJ}{[1][0][]source}{}{}{}}
\makeatother
\title{Underlying sets of limits of decent algebraic spaces}
\author{The Stacks Project Authors}
\date{stacks-project:spaces-limits.tex:local-snapshot-20261012; GNU FDL 1.2 or later}
\begin{document}
\maketitle
\noindent\textit{Editorial note.} Counted span is the author's own proof environment, source lines 1101-1153, following the statement at 1092-1099.\par\medskip
\noindent\textit{Editorial note.} No load-bearing omissions; the closing parenthesis on quasi-separated / locally separated spaces being decent is transcribed.\par\medskip
\noindent\textit{Allowed standard prerequisites (editorial).}\begin{itemize}
\item Presentation of the limit space. lemma-directed-inverse-system-has-limit (same chapter, cited at \textasciitilde{}1104-1105 and \textasciitilde{}1114-1115). X = U/R with U = lim U\_i and R = lim R\_i; also f\_0\textasciicircum{}\{-1\}W\_0 = lim f\_\{i0\}\textasciicircum{}\{-1\}W\_0.
\item Topology of limits of schemes. Limits, Lemma limits-lemma-topology-limit (cross-chapter, cited at \textasciitilde{}1117-1118). |U| = lim |U\_i| and |R| = lim |R\_i| as sets.
\item Nonemptiness of directed limits of finite nonempty sets. Categories, Lemma categories-lemma-nonempty-limit (cross-chapter, cited at \textasciitilde{}1125-1126 and \textasciitilde{}1140-1141). lim S\_i and lim T\_i are nonempty.
\item Decent spaces have finite fibres under etale covers. Decent Spaces, Definition decent-spaces-definition-very-reasonable (cross-chapter, cited at \textasciitilde{}1122-1124 and \textasciitilde{}1136-1138). S\_i and T\_i are finite nonempty sets.
\end{itemize}\medskip
\ifdefined\hypertarget\hypertarget{sourcepointer-1}{}\fi
\noindent\textit{Source reference.} Original reference 086J: focused-seeds/geometry/S9-spaces-limits-inverse-limit-sets/ref-tags.txt, source lines 5--5. Navigation pointer only: the counted proof cites this statement by \textbackslash{}ref as a standing input; the deduction the counted proof has to make is carried out inside the counted span, and no part of that deduction is delegated to this pointer.\par\medskip
\ifdefined\hypertarget\hypertarget{sourcepointer-2}{}\fi
\noindent\textit{Source reference.} Original reference 03I8: focused-seeds/geometry/S9-spaces-limits-inverse-limit-sets/ref-tags.txt, source lines 6--6. Navigation pointer only: the counted proof cites this statement by \textbackslash{}ref as a standing input; the deduction the counted proof has to make is carried out inside the counted span, and no part of that deduction is delegated to this pointer.\par\medskip
\ifdefined\hypertarget\hypertarget{sourcepointer-3}{}\fi
\noindent\textit{Source reference.} Original reference 088J: focused-seeds/geometry/S9-spaces-limits-inverse-limit-sets/ref-tags.txt, source lines 7--7. Navigation pointer only: the counted proof cites this statement by \textbackslash{}ref as a standing input; the deduction the counted proof has to make is carried out inside the counted span, and no part of that deduction is delegated to this pointer.\par\medskip
\ifdefined\hypertarget\hypertarget{sourcepointer-4}{}\fi
\noindent\textit{Source reference.} Original reference 03I7: focused-seeds/geometry/S9-spaces-limits-inverse-limit-sets/ref-tags.txt, source lines 8--8. Navigation pointer only: the counted proof cites this statement by \textbackslash{}ref as a standing input; the deduction the counted proof has to make is carried out inside the counted span, and no part of that deduction is delegated to this pointer.\par\medskip
\ifdefined\hypertarget\hypertarget{sourcepointer-5}{}\fi
\noindent\textit{Source reference.} Original lemma 07SF: sources/stacks-spaces-limits/pinned-original-source.tex, source lines 878--878. Navigation pointer only: the counted proof cites this statement by \textbackslash{}ref as a standing input; the deduction the counted proof has to make is carried out inside the counted span, and no part of that deduction is delegated to this pointer.\par\medskip
\ifdefined\hypertarget\hypertarget{sourcepointer-6}{}\fi
\noindent\textit{Source reference.} Original reference 086V: focused-seeds/geometry/S9-spaces-limits-inverse-limit-sets/ref-tags.txt, source lines 9--9. Navigation pointer only: the counted proof cites this statement by \textbackslash{}ref as a standing input; the deduction the counted proof has to make is carried out inside the counted span, and no part of that deduction is delegated to this pointer.\par\medskip
\ifdefined\hypertarget\hypertarget{sourcepointer-7}{}\fi
\noindent\textit{Source reference.} Original reference 03IJ: focused-seeds/geometry/S9-spaces-limits-inverse-limit-sets/ref-tags.txt, source lines 10--10. Navigation pointer only: the counted proof cites this statement by \textbackslash{}ref as a standing input; the deduction the counted proof has to make is carried out inside the counted span, and no part of that deduction is delegated to this pointer.\par\medskip
\begin{lemma}
\label{lemma-inverse-limit-sets}
Let $S$ be a scheme. Let $X = \lim_{i \in I} X_i$ be the limit of a directed
inverse system of algebraic spaces over $S$ with affine transition morphisms
(Lemma \ref{lemma-directed-inverse-system-has-limit}). If each $X_i$
is decent (for example quasi-separated or locally separated)
then $|X| = \lim_i |X_i|$ as sets.
\end{lemma}

\begin{proof}
There is a canonical map $|X| \to \lim |X_i|$. Choose $0 \in I$.
If $W_0 \subset X_0$ is an open subspace, then we have
$f_0^{-1}W_0 = \lim_{i \geq 0} f_{i0}^{-1}W_0$, see
Lemma \ref{lemma-directed-inverse-system-has-limit}.
Hence, if we can prove the lemma for inverse systems where $X_0$
is quasi-compact, then the lemma follows in general. Thus we may and do
assume $X_0$ is quasi-compact.

\medskip\noindent
Choose an affine scheme $U_0$ and a surjective \'etale morphism $U_0 \to X_0$.
Set $U_i = X_i \times_{X_0} U_0$ and $U = X \times_{X_0} U_0$.
Set $R_i = U_i \times_{X_i} U_i$ and $R = U \times_X U$.
Recall that $U = \lim U_i$ and $R = \lim R_i$, see proof of
Lemma \ref{lemma-directed-inverse-system-has-limit}.
Recall that $|X| = |U|/|R|$ and $|X_i| = |U_i|/|R_i|$. By
Limits, Lemma \ref{limits-lemma-topology-limit} we have
$|U| = \lim |U_i|$ and $|R| = \lim |R_i|$.

\medskip\noindent
Surjectivity of $|X| \to \lim |X_i|$. Let $(x_i) \in \lim |X_i|$. Denote
$S_i \subset |U_i|$ the inverse image of $x_i$. This is a finite nonempty
set by the definition of decent spaces
(Decent Spaces, Definition \ref{decent-spaces-definition-very-reasonable}).
Hence $\lim S_i$ is nonempty, see
Categories, Lemma \ref{categories-lemma-nonempty-limit}.
Let $(u_i) \in \lim S_i \subset \lim |U_i|$. By the above this determines
a point $u \in |U|$ which maps to an $x \in |X|$ mapping to the given
element $(x_i)$ of $\lim |X_i|$.

\medskip\noindent
Injectivity of $|X| \to \lim |X_i|$. Suppose that $x, x' \in |X|$
map to the same point of $\lim |X_i|$. Choose lifts $u, u' \in |U|$
and denote $u_i, u'_i \in |U_i|$ the images.
For each $i$ let $T_i \subset |R_i|$ be the set of points mapping
to $(u_i, u'_i) \in |U_i| \times |U_i|$. This is a finite
set by the definition of decent spaces
(Decent Spaces, Definition \ref{decent-spaces-definition-very-reasonable}).
Moreover $T_i$ is nonempty as we've assumed that $x$ and $x'$ map to the
same point of $X_i$. Hence $\lim T_i$ is nonempty, see
Categories, Lemma \ref{categories-lemma-nonempty-limit}.
As before let $r \in |R| = \lim |R_i|$ be a point corresponding to an
element of $\lim T_i$. Then $r$ maps to $(u, u')$ in $|U| \times |U|$
by construction and we see that $x = x'$ in $|X|$ as desired.

\medskip\noindent
Parenthetical statement: A quasi-separated algebraic space is decent, see
Decent Spaces, Section \ref{decent-spaces-section-reasonable-decent}
(the key observation to this is Properties of Spaces, Lemma
\ref{spaces-properties-lemma-finite-fibres-presentation}).
A locally separated algebraic space is decent by
Decent Spaces, Lemma \ref{decent-spaces-lemma-locally-separated-decent}.
\end{proof}

\end{document}
